\chapter{OAuth 2.1 Authorization Server}
\label{ch:oauth-as}

Chapter~\ref{ch:oauth} covered the half of OAuth that protects an API --- validating tokens issued elsewhere.  This chapter covers the other half: the authorization server (AS) that \emph{issues} those tokens.

Most Kiss applications that ship an MCP server eventually need both halves.  In the typical deployment, a single Kiss instance is the AS, the resource server, and the MCP server: the same JVM authenticates the user, runs the consent UI, issues an access token, and validates that token on the next MCP call.  Splitting the AS into a separate service is straightforward (point \texttt{OAuthAuthorizationServer} at the external AS, leave \texttt{OAuthAsEnabled=false}), but the collocated case is what the framework optimizes for.

The AS is deliberately minimal but spec-compliant.  It implements the parts of OAuth 2.1 the MCP specification requires --- authorization-code flow with PKCE, refresh-token rotation, dynamic client registration, the metadata and JWKS endpoints --- and stops there.  Section~\ref{sec:as-omissions} catalogues what is left out and why.

This chapter and the previous one cover Kiss as a server of OAuth.  The third role --- Kiss as a \emph{client} of a remote OAuth server, obtaining tokens to make its own authorized outbound calls --- is the subject of Chapter~\ref{ch:oauth-client}.  The two are independent: a single app may issue tokens, validate tokens, and consume a partner's tokens all at once, or use any subset.

\section{Configuration}
\label{sec:as-config}

All settings live in \texttt{src/main/back\-end/applic\-ation.ini}.  Setting \texttt{OAuth\-As\-Enabled=true} activates the AS; everything else has a sensible default.  In the common case where the AS and the resource server are the same Kiss app, only one URL setting is needed --- and it lives in the resource-server section, not here:

\begin{lstlisting}[style=ini]
# Required.
OAuthAsEnabled = true

# Optional. Canonical URL of this AS. Goes into the iss claim
# of every issued token and forms the base of the metadata
# document. Defaults to OAuthAuthorizationServer (from the
# resource-server config in the preceding section of
# application.ini) when not set --- which is the right value
# when the AS and the RS are the same Kiss app. Set this
# explicitly only when the AS publishes itself under a URL
# different from the RS-side OAuthAuthorizationServer value.
# OAuthAsIssuer = https://myapp.example.com

# Optional defaults shown
# (OAuthAsSqliteFile is relative to applicationPath, or an
# absolute path)
# OAuthAsSqliteFile             = oauth.db
# OAuthAccessTokenTtlSeconds    = 3600     # 1 hour
# OAuthRefreshTokenTtlSeconds   = 2592000  # 30 days
# OAuthAuthCodeTtlSeconds       = 60       # 1 minute
# OAuthPruneIntervalSeconds     = 900      # 15 minutes
# OAuthAllowDynamicRegistration = true     # required for MCP
# OAuthSessionTtlSeconds        = 1800     # AS login cookie
# OAuthKeyId                    = kiss-key-1  # JWT kid
\end{lstlisting}

When the AS is disabled, every endpoint returns 404 --- the same response a server with no AS would give, signaling to clients that this resource is not paired with an AS at this URL.

\begin{tipbox}
The application path (\texttt{MainServlet.getApplicationPath()}) is already the backend directory --- in dev mode \texttt{src/main/backend/}, in a deployed WAR \texttt{WEB-INF/backend/}.  So the default \texttt{OAuthAsSqliteFile=oauth.db} colocates the file with \texttt{application.ini}.  You generally do not need to change this.

\texttt{OAuthAsSqliteFile} may also be set to an \emph{absolute} path (e.g.\ \texttt{/var/lib/myapp/oauth.db}), which is used verbatim.  This is useful when you want the runtime state file to live outside the deployed webapp tree so a WAR redeploy cannot touch it --- the signing keys, registered clients, and refresh tokens then survive a redeploy unchanged.

The SQLite file is created and its schema initialized automatically on first startup; no setup tool is needed, and the database is completely independent of the application's main database (its own file, its own JDBC connection, no shared schema or connection pool).
\end{tipbox}

\section{Required App Code: UserAuthenticator}
\label{sec:as-user-auth}

The AS will refuse to issue tokens until the application provides a {\nopagebreak}\texttt{UserAuthenticator}\index{UserAuthenticator}.  This is the one piece of the AS that cannot be generic: the credential check inevitably knows your user table schema, your LDAP layout, or your IdP's API.

The interface has one method:

\begin{lstlisting}[style=java]
public interface UserAuthenticator {
    AuthenticatedUser authenticate(String username,
        String password);
}
\end{lstlisting}

\noindent
Return an \texttt{AuthenticatedUser}\index{AuthenticatedUser} on success --- carrying the user's subject identifier and an optional bag of extra claims that the framework will copy onto every access token issued for this user.  Return \texttt{null} on any kind of failure: bad credentials, locked account, user store unreachable.  Never throw --- the framework cannot distinguish ``wrong password'' from ``LDAP is down'' if you throw, and the user-facing render should not leak that distinction anyway.  Log the underlying cause inside your implementation.

For Kiss applications that already have a \texttt{Login.groovy} (Chapter~\ref{ch:authentication}), a ten-line wrapper is sufficient:

\begin{lstlisting}[style=groovy]
// src/main/backend/MyUserAuthenticator.groovy
package com.mycompany

import org.kissweb.oauth.as.AuthenticatedUser
import org.kissweb.oauth.as.UserAuthenticator
import org.kissweb.json.JSONObject

class MyUserAuthenticator implements UserAuthenticator {
    AuthenticatedUser authenticate(String username,
        String password) {
        def row = Login.verifyPassword(username, password)
        if (!row)
            return null
        def extras = new JSONObject()
        extras.put("email", row.email)
        extras.put("name",  row.full_name)
        return new AuthenticatedUser(String.valueOf(
            row.user_id), extras)
    }
}
\end{lstlisting}

\noindent
Register it in \texttt{KissInit.groovy}, after \texttt{MainServlet.read\-IniFile(...)}:

\begin{lstlisting}[style=groovy]
AsExtensions.setUserAuthenticator(
    new com.mycompany.MyUserAuthenticator())
\end{lstlisting}

\noindent
A few things worth being deliberate about:

\begin{description}[style=nextline]
  \item[Use a stable subject.]  The subject (\texttt{sub} claim) is the primary identifier every downstream component uses to recognize this user --- audit logs, resource-server authorization decisions, custom claims joined to other databases.  Prefer the database primary key over a username that can be renamed.
  \item[Be deliberate about extra claims.]  Anything in \texttt{extraClaims} ends up signed into every access token, visible to whatever resource server consumes the token.  Email is usually fine; password hashes and internal IDs other systems should not see, are not.
  \item[Implementations live in app code.]  Either \texttt{src/main/backend/} (Groovy or Java, hot-reloadable) or \texttt{src/main/precompiled/} (Java with package structure, restart required after change).
\end{description}

\section{Optional App Code: ConsentProvider and HTML Overrides}
\label{sec:as-consent}

\subsection{ConsentProvider}

\texttt{ConsentProvider}\index{ConsentProvider} supplies the human-readable text the consent page shows for each scope a client is requesting.  The framework knows scope \emph{names} like \texttt{mcp:write} but has no idea what they mean in your application's domain.  The default implementation returns the scope name verbatim --- functional but unfriendly:

\begin{lstlisting}[style=java]
public interface ConsentProvider {
    String describeScope(String scope);
    default String getDisplayName() { ... }
}
\end{lstlisting}

\noindent
Provide your own to turn cryptic scope names into sentences a user can understand:

\begin{lstlisting}[style=groovy]
AsExtensions.setConsentProvider(new ConsentProvider() {
    String describeScope(String s) {
        switch (s) {
            case "mcp:read":
                return "Read your orders, customers, " +
                    "and reports."
            case "mcp:write":
                return "Create, update, and delete records " +
                    "on your behalf."
            default:
                return s
        }
    }
    String getDisplayName() { return "Acme Corp" }
})
\end{lstlisting}

\subsection{HTML Overrides}

The login and consent pages have minimal built-in HTML that works out of the box.  To restyle or rebrand, drop replacement files in \texttt{src/main/front\-end/oauth/}:

\begin{itemize}
  \item \texttt{src/main/frontend/oauth/login.html}
  \item \texttt{src/main/frontend/oauth/consent.html}
\end{itemize}

\noindent
Both use simple \texttt{\{placeholder\}} substitution.  The supported placeholders:

\begin{description}[style=nextline]
  \item[\texttt{\{asName\}}]  \texttt{ConsentProvider.getDisplayName()}.  Available in login and consent.
  \item[\texttt{\{clientName\}}]  The DCR \texttt{client\_name} (or the client id if no name was registered).  Available in login and consent.
  \item[\texttt{\{flowId\}}]  Opaque flow id assigned to this in-flight authorization request.  Must round-trip in the form as a hidden field named \texttt{flow}.  Available in login and consent.
  \item[\texttt{\{userSubject\}}]  Subject of the user who just logged in.  Available in consent only.
  \item[\texttt{\{scopeList\}}]  Pre-rendered \texttt{<li>} items for each requested scope, with the description from \texttt{ConsentProvider}.  Available in consent only.
  \item[\texttt{\{errorBlock\}}]  An error \texttt{<div>}, or empty when there is nothing to show.  Available in login and consent.
\end{description}

Both forms must POST back to \texttt{/oauth/authorize} with a hidden \texttt{stage=login} or \texttt{stage=consent} and the \texttt{flow} id.  The defaults in \texttt{AuthorizationServlet.java} are the working reference --- copy them and rework the styling.

\begin{notebox}
Templates are loaded once at first use and cached for the lifetime of the JVM.  Edits to your template files take effect after a server restart.  Hot reload does not apply --- the servlet container does not re-scan template files between requests.
\end{notebox}

\section{The Five Endpoints}
\label{sec:as-endpoints}

All five are auto-registered via \texttt{@WebServlet} annotations; no \texttt{web.xml} entries required.

\begin{description}[style=nextline]
  \item[\texttt{GET /.well-known/oauth-authorization-server}]  RFC 8414 metadata.  Returns the issuer URL, the endpoint URLs, the supported grant types, response types, PKCE methods, token endpoint auth methods, and signing algorithm.  Cached one hour at the client.  Returns 404 when the AS is disabled.
  \item[\texttt{GET /oauth/jwks}]  A single-key JWKS document carrying the AS's current RSA-2048 public key.  \texttt{kid} matches the \texttt{kid} on every issued JWT, so a resource server can fetch this once and validate signatures locally without further round-trips.  Cached five minutes.
  \item[\texttt{GET /oauth/authorize}]  Authorization endpoint.  Validates the request, then renders the login page (no session) or the consent page (session present).  On approval, redirects the user-agent back to the client's \texttt{redirect\_uri} with \texttt{?code=...\&state=...}; on denial, with \texttt{?error=access\_denied\&...}.  Errors that prevent the redirect (unknown client, redirect URI mismatch) render as 400 HTML --- the framework refuses to redirect to a URI not registered for the client.
  \item[\texttt{POST /oauth/token}]  Token endpoint.  Handles \texttt{authorization\_code} and \texttt{refresh\_token} grants.  All responses include \texttt{Cache-Control: no-store}.
  \item[\texttt{POST /oauth/register}]  RFC 7591 dynamic client registration.  Anyone can POST; the response carries a freshly-generated \texttt{client\_id} and, for confidential clients, a \texttt{client\_secret}.  Disable via \texttt{OAuthAllowDynamic\-Registration=false} if your deployment does not need open registration.
\end{description}

\section{The Authorization Code Flow Step by Step}
\label{sec:as-flow}

The flow as it plays out across browser, AS, and client:

\begin{enumerate}
  \item Client begins the dance: builds an authorization URL and either redirects the user-agent there (web client) or opens it in a browser (native / desktop / MCP client):

\begin{lstlisting}[style=shell]
GET /oauth/authorize
    ?response_type=code
    &client_id=<id>
    &redirect_uri=<exact-match-of-registered>
    &scope=mcp:read mcp:write
    &state=<csrf nonce>
    &code_challenge=<S256 hash of verifier>
    &code_challenge_method=S256
    # optional RFC 8707
    &resource=https://myapp.example.com
\end{lstlisting}

  \item AS validates everything.  Unknown \texttt{client\_id} or mismatched \texttt{redirect\_uri} render as 400 HTML and stop here.  Other errors (unknown scope, missing PKCE) redirect back to the client with \texttt{?error=...}.
  \item If the user has no session, the AS renders the login page.  The form's hidden \texttt{flow=<id>} keys an in-memory record of the pending request.
  \item User submits credentials.  The AS calls \texttt{UserAuthenticator.authen\-ticate()}.  On failure, re-render the login page with an error.  On success, set an \texttt{oauth\_session} cookie (HttpOnly, Secure when over TLS, scoped to \texttt{/oauth/}, 30-minute TTL by default) and render the consent page.
  \item User clicks Allow.  The AS generates an opaque authorization code (60-second TTL, single-use), stores it in memory along with the user, scopes, audience, redirect URI, PKCE challenge, and any extra claims from the authenticator, then 302-redirects to:

\begin{lstlisting}[style=shell]
HTTP 302
Location: <redirect_uri>?code=<code>&state=<state>
\end{lstlisting}

  \item Client receives the redirect, extracts the code, and POSTs to the token endpoint:

\begin{lstlisting}[style=shell]
POST /oauth/token
Content-Type: application/x-www-form-urlencoded

grant_type=authorization_code
&code=<code>
&redirect_uri=<must match step 1 exactly>
&client_id=<id>
&code_verifier=<PKCE verifier>
\end{lstlisting}

  \item AS verifies: client exists, code is valid and not consumed, redirect URI matches, PKCE \texttt{SHA256(code\_verifier) == code\_challenge}.  Then it mints an RS256 access token (RFC 9068 \texttt{at+jwt} profile) carrying \texttt{iss / sub / aud / exp / iat / nbf / jti / scope / client\_id} plus any extra claims, generates an opaque refresh token, stores the refresh token, and returns:

\begin{lstlisting}[style=shell]
HTTP 200 OK
Cache-Control: no-store
Content-Type: application/json

{
  "access_token":  "eyJ...",
  "token_type":    "Bearer",
  "expires_in":    3600,
  "refresh_token": "...",
  "scope":         "mcp:read mcp:write"
}
\end{lstlisting}

  \item Client uses the access token against the resource server (Chapter~\ref{ch:oauth}) until it expires, then refreshes (Section~\ref{sec:as-rotation}).
\end{enumerate}

Subsequent \texttt{/authorize} requests from the same browser within the session window skip the login step --- the \texttt{oauth\_session} cookie identifies the user, and only consent is shown.  Logout is implicit: the cookie expires (default 30 minutes) or the user clears it.

\section{Dynamic Client Registration}
\label{sec:as-dcr}

DCR is enabled by default because the MCP specification requires it.  Anyone reachable from the AS can POST a client metadata document to \texttt{/oauth/register} and receive back a freshly-generated client id (and, for confidential clients, a one-time-displayed secret).  Persisted to the AS's SQLite database (\texttt{oauth.db}).

\begin{lstlisting}[style=shell]
curl -X POST http://localhost:8001/oauth/register \
  -H "Content-Type: application/json" \
  -d '{
    "redirect_uris": ["http://localhost:54545/cb"],
    "client_name": "Claude Desktop",
    "scope": "mcp:read mcp:write",
    "token_endpoint_auth_method": "none"
  }'
\end{lstlisting}

\noindent
\texttt{token\_endpoint\_auth\_method=none} produces a \emph{public client} (no secret).  Public clients authenticate at the token endpoint via PKCE alone; this is the right shape for native and desktop MCP clients that cannot keep a secret.

Any other supported value (\texttt{client\_secret\_basic} or \texttt{client\_secret\_post}, the default) produces a \emph{confidential client}.  The server generates a 256-bit secret, returns it in the response exactly once, and stores only its SHA-256 hash.  There is no recovery if the client loses the secret --- it must re-register.

For production where DCR is not appropriate, set \texttt{OAuthAllowDynamic\-Registration=false} and provision clients programmatically from startup code:

\begin{lstlisting}[style=groovy]
import org.kissweb.oauth.as.*

ClientStore.get().register(new RegisteredClient(
    "claude-prod",
    // pre-computed SHA-256 hex of secret
    "sha256:abc123...",
    "Claude Production",
    ["https://prod.example/cb"],
    ["mcp:read", "mcp:write"] as Set,
    ["authorization_code", "refresh_token"] as Set,
    System.currentTimeMillis() / 1000))
\end{lstlisting}

There is no RFC 7592 client-management endpoint in v1 --- to update or delete a registered client, use \texttt{ClientStore.remove(clientId)} from app code.

\section{Refresh-Token Rotation and Reuse Detection}
\label{sec:as-rotation}

OAuth 2.1 mandates refresh-token rotation for the \texttt{refresh\_token} grant, and Kiss implements it.  Every successful refresh:

\begin{enumerate}
  \item Issues a brand-new refresh token in the same \emph{family}.  The family id is a UUID assigned when the original authorization-code grant happened; every rotation in the same logical session keeps the same family id.
  \item Marks the old refresh token as rotated, with its \texttt{rotated\_to\_jti} pointing at the successor.
  \item Returns the new pair (access + refresh) to the client.
\end{enumerate}

If a \emph{rotated} refresh token is ever presented again --- which an honest client never does because it has already upgraded to the successor --- the AS treats this as evidence of token theft and \textbf{revokes the entire family}.  Every refresh token in the family is deleted from the AS database; the next refresh attempt by anyone in the family (including the legitimate client) returns \texttt{invalid\_grant}.  The legitimate user has to log in again.

\begin{notebox}
Rotation does not prevent the initial theft.  It bounds the damage: the attacker cannot keep refreshing indefinitely without the legitimate client noticing the next time \emph{it} tries to refresh, at which point both attacker and legitimate client lose access.
\end{notebox}

This is the OAuth 2.1 stolen-refresh-token detection mechanism described in RFC 6749bis \S{}6.1.  Implementation-wise, rotation issues two writes (mark the old token rotated; insert the successor) inside one SQLite transaction so the on-disk state is never observed mid-rotation; family revocation is a single indexed \texttt{DELETE}.  Both complete in milliseconds at MCP-scale traffic --- see Section~\ref{sec:as-persistence}.

\section{Persistence: The oauth.db Database}
\label{sec:as-persistence}

All long-lived AS state lives in a private SQLite database --- a single file (\texttt{oauth.db} by default), independent of the application's main database.  The AS opens its own JDBC connection against this file, shares no connection pool with the rest of the application, and requires no operator setup beyond pointing at a file path.  The schema is materialized via \texttt{CREATE TABLE IF NOT EXISTS} on every open, so a fresh install creates the file automatically.

\begin{lstlisting}[style=sql]
CREATE TABLE oauth_keys (
    kid          TEXT PRIMARY KEY,
    private_key  TEXT NOT NULL,           -- base64 PKCS#8
    public_key   TEXT NOT NULL,           -- base64 X.509
    created_at   INTEGER NOT NULL
);

CREATE TABLE oauth_clients (
    client_id            TEXT PRIMARY KEY,
    -- null for public clients
    client_secret_hash   TEXT,
    client_name          TEXT,
    redirect_uris        TEXT,            -- space-joined
    allowed_scopes       TEXT,            -- space-joined
    allowed_grant_types  TEXT,            -- space-joined
    created_at           INTEGER NOT NULL
);

CREATE TABLE oauth_refresh_tokens (
    jti                TEXT PRIMARY KEY,
    family_id          TEXT NOT NULL,
    client_id          TEXT NOT NULL,
    user_sub           TEXT,
    user_extra_claims  TEXT,              -- base64 JSON
    scopes             TEXT,              -- space-joined
    audience           TEXT,
    created_at         INTEGER NOT NULL,
    expires_at         INTEGER NOT NULL,
    -- null for active, set for rotated
    rotated_to_jti     TEXT
);
CREATE INDEX idx_oauth_refresh_family
    ON oauth_refresh_tokens (family_id);
CREATE INDEX idx_oauth_refresh_expires
    ON oauth_refresh_tokens (expires_at);

CREATE TABLE oauth_meta (
    meta_key    TEXT PRIMARY KEY,
    -- holds schema_version etc.
    meta_value  TEXT NOT NULL
);
\end{lstlisting}

\noindent
Authorization codes (60-second TTL) live in memory only --- losing one on restart costs the user one click of the Allow button, not worth a write.

\subsection{Transactions}

Every mutating operation commits as a normal SQLite transaction.  Refresh-token rotation in particular issues both writes (mark the old token; insert the successor) inside one transaction, so the on-disk state is never observed mid-rotation.  A JVM crash before commit leaves the previous state intact --- no half-written rows, no all-keys-lost catastrophe.

\subsection{Lazy Pruning}

There is no cron job.  Every OAuth endpoint call checks an in-memory \texttt{last\-Pruned\-At} timestamp; if \texttt{OAuth\-Prune\-Interval\-Seconds} (default 15 minutes) has elapsed since the last prune, the call issues a single \texttt{DELETE \dots\ WHERE expires\_at <= ?} for refresh tokens and drops expired in-memory codes.  The \texttt{expires\_at} column is indexed so the sweep is cheap regardless of table size.  The throttle timestamp is in memory only --- a restart causes the first call to re-prune, which is exactly what we want.

\subsection{Concurrency Model}

Single-process only.  SQLite serializes writers internally, and the Kiss \texttt{Connection} is shared across threads under a JVM-wide monitor on \texttt{OAuthSqliteStore}.  Two Kiss JVMs writing to the same \texttt{oauth.db} will fail with SQLite's standard \texttt{database is locked} error rather than silently corrupting state.  For high availability, run a single AS instance behind a load balancer, or swap in alternative implementations of \texttt{ClientStore}, \texttt{RefreshTokenStore}, and \texttt{KeyManager} that talk to a shared server-class database.

\subsection{Scale}

Per-row overhead is well under a kilobyte.  Tens of thousands of active refresh tokens fit comfortably in a single-MB file with sub-millisecond point lookups (\texttt{jti}, \texttt{family\_id}, and \texttt{expires\_at} are all indexed).  Even after years of accumulated registered clients, the database remains small enough to ship in backups and trivially small to inspect.

\subsection{Inspection}

Operators with the \texttt{sqlite3} CLI can inspect the file in place:

\begin{lstlisting}[style=shell]
sqlite3 /path/to/oauth.db '.tables'
sqlite3 /path/to/oauth.db \
  'SELECT client_id, client_name FROM oauth_clients;'
\end{lstlisting}

\noindent
Take a copy first if the AS is running.

\subsection{Build-Time Safety}

\texttt{./bld build} copies \texttt{src/main/backend/} into the deployed WAR, but the framework's build script excludes \texttt{oauth.db} (and SQLite's transactional sidecar files \texttt{oauth.db-journal}, \texttt{-wal}, \texttt{-shm}) from the copy.  An accidental \texttt{oauth.db} in the source tree will never overwrite the runtime file in the deployed app.  Also add the same set of names to \texttt{.gitignore} --- runtime state should never be checked in.

For deployments where even \texttt{WEB-INF/backend/} is rewritten on redeploy (e.g.\ a full WAR replacement rather than an in-place update), set \texttt{OAuthAsSqliteFile} to an absolute path outside the webapp directory --- the file is then completely untouched by the deploy and the AS retains its signing keys, registered clients, and refresh tokens across upgrades.

\section{Security Considerations}
\label{sec:as-security}

\begin{description}[style=nextline]
  \item[HTTPS in production.]  The session cookie carries the user's identity for the duration of the consent flow; access and refresh tokens are bearer credentials.  TLS termination is the deployer's responsibility --- Kiss does not enforce it.  Cookies are automatically marked \texttt{Secure} when the request is HTTPS.
  \item[PKCE is mandatory.]  S256 only; \texttt{plain} is rejected even when explicitly requested.  Public clients depend entirely on PKCE for authentication.
  \item[Exact redirect-URI matching.]  Per OAuth 2.1 \S{}4.1.2.1, no wildcards, no scheme/host normalization, no fragments.  The \texttt{redirect\_uri} on \texttt{/authorize} must equal one of the client's registered URIs byte-for-byte, and the same value must come back on \texttt{/token}.  The framework refuses to redirect to a non-registered URI even when the client says to.
  \item[Refresh-token rotation is mandatory.]  See Section~\ref{sec:as-rotation}.  Without rotation, a leaked refresh token is good for 30 days with no detection mechanism.
  \item[Audience binding.]  If the client supplied a \texttt{resource} parameter on \texttt{/authorize}, the issued tokens carry it as the \texttt{aud} claim and the audience is enforced on \texttt{/token} refresh.  Otherwise, tokens are bound to the AS issuer URL --- any resource server that trusts this AS will accept them.
  \item[Client secret storage.]  Only SHA-256 of the secret is stored.  The full secret is returned to the client exactly once during DCR; there is no recovery if it's lost.
  \item[\texttt{alg=none} cannot be issued.]  \texttt{TokenIssuer} signs with \texttt{SHA256withRSA} directly; there is no algorithm-from-header dispatch on the issuing side.
  \item[Authorization codes are single-use.]  \texttt{AuthorizationCodeStore.consume()} returns the code and removes it from the store in one synchronized operation.  A replay attempt finds nothing.
  \item[Client secrets sent in basic auth are constant-time compared] against the stored hash to avoid leaking partial matches via timing.
  \item[No initial-access-token gate on DCR.]  Anyone reachable from the AS can register a client.  This is desirable for MCP (clients you've never met need to onboard) but means the registry can grow unboundedly under attack.  Disable DCR with \texttt{OAuthAllowDynamicRegistration=false} if open registration is not acceptable for your deployment.
\end{description}

\section{Components}
\label{sec:as-components}

All classes live in \texttt{org.kiss\-web.oauth.as}\index{org.kissweb.oauth.as}.  Application developers do not write or modify any of them --- only the two interfaces \texttt{UserAuthenticator} and \texttt{ConsentProvider} (in app code), plus \texttt{appli\-cation.ini} keys.

\begin{description}[style=nextline]
  \item[\texttt{AuthorizationServerConfig}\index{AuthorizationServerConfig}]  Reads AS config from \texttt{application.ini}.  Lazy singleton.
  \item[\texttt{OAuthSqliteStore}\index{OAuthSqliteStore}]  Opens and initializes the schema of the AS's private SQLite database, provides the shared JDBC \texttt{Connection} to the other stores, and holds the in-memory prune-throttle timestamp.
  \item[\texttt{KeyManager}\index{KeyManager}]  Generates RSA-2048 keypair on first start, persists as base64 PKCS\#8 / X.509 in the \texttt{oauth\_keys} table, exposes private key for signing and public JWK for publishing.
  \item[\texttt{ClientStore}\index{ClientStore}]  CRUD for registered clients in the \texttt{oauth\_clients} table.
  \item[\texttt{RefreshTokenStore}\index{RefreshTokenStore}]  CRUD for refresh tokens in the \texttt{oauth\_refresh\_tokens} table.  Implements rotation and family revocation via indexed SQL.
  \item[\texttt{AuthorizationCodeStore}\index{AuthorizationCodeStore}]  In-memory single-use codes with TTL expiry.
  \item[\texttt{TokenIssuer}\index{TokenIssuer}]  Builds and signs RS256 JWTs per RFC 9068; generates random opaque codes and refresh-token values.
  \item[\texttt{PkceValidator}\index{PkceValidator}]  RFC 7636 S256 verifier.
  \item[\texttt{AsExtensions}\index{AsExtensions}]  Static registry where the application installs its \texttt{User\-Authenticator} and \texttt{ConsentProvider}.
  \item[\texttt{UserAuthenticator}\index{UserAuthenticator} \emph{(interface)}]  Application implements: verify credentials, return \texttt{AuthenticatedUser} or null.
  \item[\texttt{ConsentProvider}\index{ConsentProvider} \emph{(interface)}]  Application implements: per-scope human description, AS display name.
  \item[\texttt{AuthenticatedUser}\index{AuthenticatedUser}]  Subject identifier plus extra claims to embed in tokens.
  \item[\texttt{RegisteredClient}\index{RegisteredClient}, \texttt{Refresh\-Token}, \texttt{Authorization\-Code}]  Immutable record types.
  \item[\texttt{MetadataServlet}\index{MetadataServlet}, \texttt{JwksPublishServlet}\index{JwksPublishServlet}]  RFC 8414 and JWKS endpoints.
  \item[\texttt{AuthorizationServlet}\index{AuthorizationServlet}]  Login, consent, code issuance.
  \item[\texttt{TokenServlet}\index{TokenServlet}]  Code exchange and refresh rotation.
  \item[\texttt{DynamicClientRegistrationServlet}\index{DynamicClientRegistrationServlet@Dynamic\-Client\-Registration\-Servlet}]  RFC 7591.
\end{description}

\section{A Worked Example}
\label{sec:as-example}

Take a Kiss application that already has a working MCP server (Chapter~\ref{ch:external}) and an existing user table.  Adding self-hosted OAuth takes three changes.

\subsection{Step 1: Turn on the AS in application.ini}

When the AS and the resource server are the same Kiss app at the same URL --- the common MCP shape --- the issuer URL appears exactly once.  \texttt{OAuthAsIssuer} defaults to \texttt{OAuthAuthorizationServer}; \texttt{OAuthResourceIdentifier} also defaults to \texttt{OAuthAuthorizationServer}; and \texttt{OAuthJwksUri} defaults to RFC 8414 discovery against the same URL (Section~\ref{sec:oauth-jwks-discovery}), which Kiss's own AS publishes (next bullet of the next section).  So the minimum config is:

\begin{lstlisting}[style=ini]
# Authorization-server config
OAuthAsEnabled = true

# Resource-server config so MCP servlets validate tokens
# issued here
OAuthAuthorizationServer = https://myapp.example.com
OAuthRequiredScopes      = mcp:read mcp:write
\end{lstlisting}

\noindent
If the AS and RS are split across different URLs --- unusual for MCP but supported --- set \texttt{OAuthAsIssuer}, \texttt{OAuthAuthorizationServer}, and \texttt{OAuthResourceIdentifier} to their respective canonical URLs.

\subsection{Step 2: Provide a UserAuthenticator}

\begin{lstlisting}[style=groovy]
// src/main/backend/MyUserAuthenticator.groovy
package com.mycompany

import org.kissweb.oauth.as.AuthenticatedUser
import org.kissweb.oauth.as.UserAuthenticator
import org.kissweb.json.JSONObject

class MyUserAuthenticator implements UserAuthenticator {
    AuthenticatedUser authenticate(String username,
        String password) {
        def row = Login.verifyPassword(username, password)
        if (!row) return null
        def extras = new JSONObject()
        extras.put("email", row.email)
        return new AuthenticatedUser(String.valueOf(
            row.user_id), extras)
    }
}
\end{lstlisting}

\subsection{Step 3: Register the authenticator and consent provider}

\begin{lstlisting}[style=groovy]
// KissInit.groovy, in init() after readIniFile(...)
AsExtensions.setUserAuthenticator(
    new com.mycompany.MyUserAuthenticator())
AsExtensions.setConsentProvider(new ConsentProvider() {
    String describeScope(String s) {
        switch (s) {
            case "mcp:read":  return "Read your data."
            case "mcp:write": return "Modify your data."
            default:          return s
        }
    }
    String getDisplayName() { return "Acme MCP" }
})
\end{lstlisting}

\subsection{That's It}

Rebuild and restart:

\begin{lstlisting}[style=shell]
./bld build
\end{lstlisting}

\noindent
The first MCP client to connect will:

\begin{enumerate}
  \item Discover the resource via \texttt{GET} \url{https://myapp.example.com/.well-known/oauth-protected-resource}.
  \item Discover the AS via \texttt{GET} \url{https://myapp.example.com/.well-known/oauth-authorization-server} --- the same metadata document Kiss's own resource server consults (Section~\ref{sec:oauth-jwks-discovery}) to find the JWKS, so the AS+RS pair needs no \texttt{OAuthJwksUri} setting.
  \item Register itself via \texttt{POST} \url{https://myapp.example.com/oauth/register}.
  \item Open a browser to \url{https://myapp.example.com/oauth/authorize?...} for the user to log in and approve.
  \item Exchange the code at \texttt{POST} \url{https://myapp.example.com/oauth/token}.
  \item Use the access token on MCP calls; refresh as needed.
\end{enumerate}

\noindent
No further code on your side.

\section{Troubleshooting}
\label{sec:as-troubleshooting}

\begin{description}[style=nextline]
  \item[Every \texttt{/oauth/...} returns 404.]  AS not enabled.  Set \texttt{OAuthAsEnabled = true} and restart (configuration is read once at startup).
  \item[\texttt{/oauth/authorize} returns 500 with ``No UserAuthenticator has been registered.'']  Add an \texttt{AsExtensions.setUserAuthenticator(...)} call to \texttt{KissInit.groovy}.
  \item[\texttt{/oauth/token} returns \texttt{invalid\_grant}: ``redirect\_uri does not match\dots''.]  OAuth 2.1 requires the same \texttt{redirect\_uri} on \texttt{/authorize} and on \texttt{/token}, byte-for-byte.  Make sure the client is sending an identical value at both calls (trailing slashes, port numbers, casing all matter).
  \item[\texttt{/oauth/token} returns \texttt{invalid\_grant}: ``PKCE code\_verifier does not match\dots''.]  The client either supplied a different verifier than the challenge was generated from, or built the challenge wrong.  Check that the client is computing \texttt{base64url(SHA-256(verifier))} (no padding) and submitting the original verifier on \texttt{/token}.
  \item[\texttt{/oauth/token} returns \texttt{invalid\_grant}: ``Refresh token has been rotated; family revoked''.]  The client (or someone with a copy of an old refresh token) submitted a rotated token.  All tokens in the family are now revoked; the user must re-authenticate.
  \item[Login page rejects valid credentials.]  Check \texttt{tomcat/logs/catalina.out} for exceptions thrown by your \texttt{UserAuthenticator}.  Exceptions render as ``Invalid username or password'' to the user but are logged in full.
  \item[\texttt{oauth.db} (or its \texttt{-journal}/\texttt{-wal}/\texttt{-shm} sidecars) appeared in \texttt{src/main/backend/} and got committed.]  Add the same set of names to \texttt{.gitignore} immediately.  The build excludes them from the WAR copy, but git does not know about that exclusion.
  \item[Issued tokens fail validation at the resource server with ``Token iss does not match expected''.]  \texttt{OAuthAsIssuer} (used as \texttt{iss} on issued tokens) must equal \texttt{OAuthAuthorizationServer} (used by the validator to check \texttt{iss}).  Set both to the same canonical URL.
\end{description}

\section{What's Not Implemented}
\label{sec:as-omissions}

The framework deliberately ships a minimal-but-spec-compliant AS.  Out of scope for v1:

\begin{itemize}
  \item \emph{Algorithms other than RS256.}  Only RS256 is issued (and only RS\{256,384,512\} can be validated).  ECDSA (ES256), EdDSA, HMAC are not implemented.  ES256 in particular is a localized addition to \texttt{TokenIssuer} and \texttt{Bearer\-Token\-Validator}.
  \item \emph{Key rotation tooling.}  The framework supports a single current key.  The JWKS structure can accommodate a published ``previous'' key during a rotation window, but the operational steps (generate new, keep old until last refresh expires, drop old) are not yet built.
  \item \emph{RFC 9449 DPoP} (proof-of-possession tokens).  All issued tokens are plain bearer tokens.
  \item \emph{RFC 8705 mTLS} client authentication or token binding.
  \item \emph{RFC 7592} client-management endpoint (update or delete a registered client via HTTP).  Use \texttt{ClientStore.remove()} from Java code.
  \item \emph{RFC 7009 revocation} and \emph{RFC 7662 introspection} endpoints.  Revocation can be done by deleting the refresh token (or calling \texttt{Refresh\-Token\-Store.revokeFamily()}); introspection is generally unnecessary when access tokens are JWTs the resource server can validate locally.
  \item \emph{OpenID Connect.}  No \texttt{openid} scope handling, no ID tokens, no userinfo endpoint.  The framework is OAuth 2.1 only; layer OIDC on top, or run a separate OIDC server.
  \item \emph{Multi-tenancy or realms.}  Single AS per Kiss instance.
  \item \emph{Shared or clustered AS storage.}  The AS backs onto a private SQLite file (\texttt{oauth.db}).  This is intentional --- the AS does not couple to the application's main database, and operators do not need to provision a schema.  For a multi-instance HA deployment, swap in alternative implementations of \texttt{ClientStore}, \texttt{RefreshTokenStore}, and \texttt{KeyManager} that talk to a shared server-class database; the rest of the AS code does not need to change.
\end{itemize}

\section{Further Reading}
\label{sec:as-references}

A standalone OAuth reference (covering both the resource and authorization server sides) lives at the repository root as \texttt{OAuth.md}.

Relevant RFCs:

\begin{itemize}
  \item \emph{RFC 6749} --- The OAuth 2.0 Authorization Framework
  \item \emph{RFC 7636} --- Proof Key for Code Exchange by OAuth Public Clients (PKCE)
  \item \emph{RFC 7517} --- JSON Web Key (JWK)
  \item \emph{RFC 7519} --- JSON Web Token (JWT)
  \item \emph{RFC 7591} --- OAuth 2.0 Dynamic Client Registration Protocol
  \item \emph{RFC 8414} --- OAuth 2.0 Authorization Server Metadata
  \item \emph{RFC 8707} --- Resource Indicators for OAuth 2.0
  \item \emph{RFC 9068} --- JSON Web Token (JWT) Profile for OAuth 2.0 Access Tokens
  \item \emph{draft-ietf-oauth-v2-1} (a.k.a.\ OAuth 2.1) --- the consolidation that mandates PKCE and refresh-token rotation
\end{itemize}

\noindent
The Model Context Protocol authorization specification (\url{https://modelcontextprotocol.io}) describes how MCP clients discover the AS, register dynamically, and obtain tokens; required reading for anyone building MCP clients or servers intended to interoperate with the broader MCP ecosystem.
